% Source: https://artofproblemsolving.com/community/c3h2675455_logic_question
% Stable identity: aops-topic-2675455-q01
\emph{Two weighings suffice.} Weigh bar $1$ against bar $2$.
\begin{itemize}
\item If they balance, both are genuine and the counterfeit is bar $3$ or $4$. Weigh $1$ against $3$: if they balance, bar $4$ is counterfeit; otherwise bar $3$ is.
\item If they do not balance, the counterfeit is bar $1$ or $2$. Weigh $1$ against $3$ (genuine): if they balance, bar $2$ is counterfeit; otherwise bar $1$ is.
\end{itemize}
\emph{One weighing is not enough.} A single weighing has only three possible outcomes (left pan down, balance, right pan down), but four different bars must be told apart, so some outcome is consistent with two different counterfeit bars. In detail: with unequal numbers of bars on the pans, the fuller pan always goes down (the difference is slight), which gives no information. With one bar against one bar, a balance leaves two candidates. With two bars against two, the scale always tips, and every bar is still possible (it could be heavy on the lower side or light on the higher side).

So the minimum is $\boxed{2}$ weighings.
